\documentclass[11pt]{article}
\usepackage[a4paper,margin=1in]{geometry}
\usepackage{fontspec}
\usepackage[boldfont=false]{xeCJK}
\setCJKmainfont{FandolSong-Regular.otf}
\usepackage{amsmath}
\usepackage{amssymb}
\usepackage{array}
\usepackage{booktabs}
\usepackage{graphicx}
\usepackage{adjustbox}
\usepackage{enumitem}
\setlist[itemize]{topsep=2pt,partopsep=0pt,parsep=0pt,itemsep=2pt,leftmargin=1.4em}
\setlist[enumerate]{topsep=2pt,partopsep=0pt,parsep=0pt,itemsep=2pt,leftmargin=1.6em}
\usepackage{parskip}
\usepackage{fvextra}
\fvset{breaklines=true,breakanywhere=true,fontsize=\small}
\setlength{\parindent}{0pt}

\begin{document}

\begin{center}
  {\LARGE\bfseries Government and the macroeconomy}\\[0.35em]
  {\large Igcse Economics}
\end{center}
\vspace{0.6em}

\section*{Government macroeconomic aims}

Governments try to manage the whole economy — the \textbf{macroeconomy} 宏观经济. They have five main aims:

\begin{enumerate}
\item \textbf{economic growth} 经济增长 — the economy makes more goods and services each year.

\item \textbf{full employment} 充分就业 — almost everyone who wants a job has one, so \textbf{unemployment} 失业 is low.

\item \textbf{price stability} 价格稳定 — prices rise only slowly, so \textbf{inflation} 通货膨胀 is low.

\item \textbf{balance of payments} 国际收支 stability — the money flowing in and out of the country (from trade) is roughly even.

\item \textbf{redistribution of income} 收入再分配 — sharing income more fairly, so the gap between rich and poor is not too wide.

\end{enumerate}

These aims can \textbf{conflict} 冲突 (work against each other). For example, fast growth can pull prices up, causing inflation. Cutting unemployment can do the same. So a government must balance its aims.

\section*{Fiscal policy}

\textbf{Fiscal policy} 财政政策 means the government changing two things to manage the economy:

\begin{itemize}
\item \textbf{government spending} 政府支出 — money the government spends, for example on schools, hospitals, and roads.

\item \textbf{taxation} 税收 — money the government takes from people and firms.

\end{itemize}

The government \textbf{budget} 预算 sets spending against tax. There are three outcomes:

\begin{itemize}
\item a \textbf{budget deficit} 预算赤字 — spending is more than the tax raised, so the government must borrow.

\item a \textbf{budget surplus} 预算盈余 — the tax raised is more than spending.

\item a balanced budget — the two are equal.

\end{itemize}

\subsection*{Types of tax}

Taxes are grouped in two ways. By \emph{what} is taxed:

\begin{itemize}
\item a \textbf{direct tax} 直接税 is paid straight to the government from \textbf{income} 收入 or \textbf{profit} 利润 (income tax, corporation tax).

\item an \textbf{indirect tax} 间接税 is added to the price of goods and paid when you spend (a sales tax).

\end{itemize}

By \emph{how the rate changes with income}:

\begin{itemize}
\item a \textbf{progressive tax} 累进税 takes a \emph{bigger} \% from higher incomes. This helps redistribute income.

\item a \textbf{regressive tax} 累退税 takes a \emph{smaller} \% from higher incomes, so it hits the poor harder.

\item a \textbf{proportional tax} 比例税 takes the \emph{same} \% from everyone.

\end{itemize}

\subsection*{What fiscal policy does}

To boost the economy (raise growth and jobs), the government can spend more or cut taxes. This raises total \textbf{demand} 需求 — but if demand grows too fast, it can cause inflation.

To slow the economy (fight inflation), the government can spend less or raise taxes. This lowers demand, but can raise unemployment.

\section*{Monetary policy}

\textbf{Monetary policy} 货币政策 is the \textbf{central bank} 中央银行 changing \textbf{interest rates} 利率 (and the \textbf{money supply} 货币供给) to manage the economy. The interest rate is the cost of borrowing and the reward for saving.

\begin{itemize}
\item \textbf{Lower} interest rates: borrowing is cheaper, so people and firms spend and \textbf{invest} 投资 more. Demand rises, giving more growth and jobs — but maybe more inflation.

\item \textbf{Higher} interest rates: borrowing is dearer, so spending falls. Demand falls, giving less inflation — but maybe more unemployment.

\end{itemize}

\section*{Supply-side policy}

\textbf{Supply-side policy} 供给侧政策 tries to raise the amount the economy \emph{can} produce. It improves the quantity and quality of the \textbf{factors of production} 生产要素. Methods include:

\begin{itemize}
\item \textbf{education and training} — better skills raise \textbf{productivity} 生产率.

\item \textbf{privatisation} 私有化 — selling state-owned firms to private owners, who may run them better.

\item \textbf{deregulation} 放松管制 — removing rules that hold firms back.

\item \textbf{incentives} 激励 — lower taxes on income and profit, so people work harder and firms invest more.

\end{itemize}

Supply-side policy can raise growth and jobs \emph{and} lower inflation at the same time. Its weakness is that it works slowly.

\section*{Economic growth}

Economic growth is a rise in \textbf{real GDP}.

\textbf{GDP} (\textbf{gross domestic product} 国内生产总值) is the total value of all goods and services a country produces in a year. \textbf{Real} 实际 GDP is GDP after the effect of rising prices is removed, so it shows the \emph{true} change in output.

\subsection*{The business cycle}

Real GDP does not rise smoothly. It moves through a \textbf{business cycle} 经济周期 with four stages:

\begin{itemize}
\item \textbf{boom} 繁荣 — fast growth; output and jobs are high, but prices rise quickly.

\item \textbf{recession} 衰退 — output falls and jobs are lost. (A recession is often defined as falling real GDP for six months.)

\item \textbf{slump} 萧条 — the bottom of the cycle; very low output and high unemployment.

\item \textbf{recovery} 复苏 — output starts to rise again.

\end{itemize}

\subsection*{Causes and effects}

Growth comes from more or better factors of production: investment in new machines, a bigger or more skilled workforce, and new technology.

\textbf{Good effects}: higher incomes, more jobs, better \textbf{living standards} 生活水平, and more tax for the government to spend.

\textbf{Bad effects}: prices may rise (inflation); the gap between rich and poor may grow; and there may be more pollution and use of scarce resources.

A recession brings the opposite: lower incomes, more unemployment, and business failures.

\section*{Employment and unemployment}

\subsection*{Measuring unemployment}

Unemployment is people who can work and want to work, but cannot find a job. There are two ways to measure it:

\begin{itemize}
\item the \textbf{claimant count} 失业救济申领人数 — it counts people who claim \textbf{unemployment benefit} 失业救济金. It misses those who want work but cannot claim.

\item the \textbf{labour force survey} 劳动力调查 — a survey that counts people who are not working but are looking for work and are ready to start.

\end{itemize}

The \textbf{labour force} 劳动力 is everyone who is working plus everyone who is unemployed (and looking for work).

\subsection*{Types of unemployment}

\begin{itemize}
\item \textbf{frictional unemployment} 摩擦性失业 — short-term unemployment between jobs (a worker who just left one job and is looking for the next).

\item \textbf{structural unemployment} 结构性失业 — a whole industry shrinks and workers' skills are no longer wanted (for example, coal miners after the mines close).

\item \textbf{cyclical unemployment} 周期性失业 — caused by a recession, when total demand is low so firms need fewer workers. (It is also called demand-deficient unemployment.)

\end{itemize}

\subsection*{Effects of unemployment}

\begin{itemize}
\item workers lose income and may lose their skills.

\item the country loses the output those workers could have made.

\item the government pays more in benefits and collects less in tax.

\item crime and poor health may rise.

\end{itemize}

\section*{Inflation}

Inflation is a rise in the \emph{average} level of prices over time. The opposite — falling prices — is \textbf{deflation} 通货紧缩.

\subsection*{Measuring inflation}

Inflation is measured by the \textbf{Consumer Prices Index} 消费者价格指数 (CPI). We pick a "basket" of goods a typical family buys, give each good a \textbf{weight} 权重 (a bigger weight for things people spend more on), and track how the price of the whole basket changes each year.

\subsection*{Causes}

\begin{itemize}
\item \textbf{demand-pull} 需求拉动 inflation — total demand grows faster than the economy can supply, so prices are pulled up.

\item \textbf{cost-push} 成本推动 inflation — firms' costs rise (wages, materials), so they raise prices to cover them.

\end{itemize}

\subsection*{Effects of inflation}

\begin{itemize}
\item money loses \textbf{purchasing power} 购买力 — each unit of money buys less.

\item savings become worth less.

\item firms must keep changing prices on labels and lists, which costs time and money.

\item the country's goods become dearer than other countries' goods, which hurts exports.

\end{itemize}

\subsection*{Deflation}

Falling prices sound good, but deflation can be harmful. People put off buying, hoping for lower prices later. So demand falls, firms cut output, and unemployment rises.


\end{document}
